\documentclass[11pt,a4paper]{article}

% ─── Encoding & Fonts ───
\usepackage[utf8]{inputenc}
\usepackage[T1]{fontenc}
\usepackage{titlesec}
\usepackage{mathpazo}
\usepackage[scaled=0.85]{helvet}

% ─── Layout ───
\usepackage[
  top=2.5cm, bottom=2.5cm,
  left=2.5cm, right=2.5cm
]{geometry}
\usepackage{setspace}
\onehalfspacing

% ─── Packages ───
\usepackage{cite}
\usepackage{graphicx}
\usepackage{booktabs}
\usepackage{longtable}
\usepackage{array}
\usepackage{multirow}
\usepackage{xcolor}
\usepackage{amsmath,amssymb}
\usepackage{hyperref}
\usepackage{cleveref}
\usepackage{enumitem}
\usepackage{caption}
\usepackage{subcaption}
\usepackage{titlesec}
\usepackage{fancyhdr}
\usepackage{etoolbox}
\usepackage{microtype}
\usepackage{parskip}
\usepackage{float}

% ─── Colors ───
\definecolor{linkblue}{HTML}{0B5394}
\definecolor{citegreen}{HTML}{38761D}
\definecolor{accentgray}{HTML}{4A4A4A}
\definecolor{rulegray}{HTML}{CCCCCC}

% ─── Hyperref Setup ───
\hypersetup{
  colorlinks=true,
  linkcolor=linkblue,
  citecolor=citegreen,
  urlcolor=linkblue,
  pdftitle={CRIMENET: The Global Network of Organized Crime},
  pdfauthor={Alvaro Franco Martins},
}

% ─── Section Formatting ───
\titleformat{\section}
  {\Large\bfseries\sffamily}{\thesection.}{0.5em}{}
\titleformat{\subsection}
  {\large\bfseries\sffamily}{\thesubsection}{0.5em}{}
\titleformat{\subsubsection}
  {\normalsize\bfseries\sffamily}{\thesubsubsection}{0.5em}{}

\titlespacing*{\section}{0pt}{2ex plus 1ex}{1ex plus 0.5ex}
\titlespacing*{\subsection}{0pt}{1.5ex plus 0.5ex}{0.75ex plus 0.25ex}
\titlespacing*{\subsubsection}{0pt}{1ex plus 0.5ex}{0.5ex plus 0.25ex}

% ─── Header / Footer ───
\pagestyle{fancy}
\fancyhf{}
\fancyhead[L]{\small\sffamily\textcolor{accentgray}{CRIMENET}}
\fancyfoot[C]{\small\sffamily\thepage}
\renewcommand{\headrulewidth}{0.4pt}

% ─── Caption styling ───
\captionsetup{
  font={small},
  labelfont={bf,sf},
  labelsep=period,
  margin=1cm,
}

% ─── Table column types ───
\newcolumntype{L}[1]{>{\raggedright\arraybackslash}p{#1}}
\newcolumntype{C}[1]{>{\centering\arraybackslash}p{#1}}

% ─── Graphics path ───
\graphicspath{{images/}}

% ═══════════════════════════════════════════════════════
%                     DOCUMENT
% ═══════════════════════════════════════════════════════
\begin{document}

% ─── Title Block ───
\begin{center}
  {\LARGE\bfseries\sffamily
  CRIMENET: The global network of organized crime}\\[1.5em]
  {\large \href{https://alvarofrancomartins.com/}{Alvaro Franco Martins, PhD}}\\[0.75em]
  {\small\textcolor{accentgray}{\today}}
\end{center}

% ─── Abstract ───
\begin{abstract}
\noindent
CRIMENET is an open-source database mapping the alliances and rivalries between criminal organizations worldwide. An LLM pipeline extracted 1,890 organizations and 3,354 relationships (2,293 alliances, 1,061 rivalries) from 771 Wikipedia articles. The alliance and rivalry networks are analyzed separately. finish by writing all the results

\medskip
\noindent\textbf{Keywords:} criminal networks, network science, signed networks, organized crime, open-source intelligence, large language models
\end{abstract}

% ═══════════════════════════════════════════════════════
\section{Introduction}
% ═══════════════════════════════════════════════════════

Organized crime has been mostly studied separately. Relationships between organizations have not been mapped at scale. Here I present CRIMENET, an open-source database of alliances and rivalries between 1{,}890 criminal organizations worldwide, extracted from 771 Wikipedia articles with an LLM pipeline. 

The complete pipeline (source URLs, extraction scripts, cleanup logic, deduplication dictionaries, and analysis code) is open-source on GitHub.\footnote{\url{https://github.com/alvarofrancomartins/CRIMENET}} An interactive D3.js visualization of the full network, with alliance/rivalry filtering and organization search, is deployed on my website.\footnote{\url{https://www.alvarofrancomartins.com/crimenet}}

This report is structured as follows: Section~\ref{sec:data} describes the data. Section~\ref{sec:network} analyzes the networks. Section~\ref{sec:limitations} lists current limitations. Section~\ref{sec:conclusion} concludes.

% ═══════════════════════════════════════════════════════
\section{Data}\label{sec:data}
% ═══════════════════════════════════════════════════════

\subsection{Source collection}
I manually compiled a list of 771 Wikipedia article URLs about criminal organizations, aiming for broad coverage across cartels, mafias, gangs, triads, motorcycle clubs, militias, and terrorist groups. 

Wikipedia articles are edited continuously, and the generic URL always resolves to the latest revision. To make the dataset reproducible, I recorded a permanent revision URL\footnote{Example: \url{https://en.wikipedia.org/w/index.php?title=\%27Ndrangheta&oldid=1341668053}} for each article rather than the generic one. The text used for extraction was fetched at this exact revision, so every organization and relationship in the dataset traces back to the precise text it was extracted from, regardless of later edits. 

For each article, I made two calls to the MediaWiki Action API: one for the clean plain-text body and one for the rendered HTML. Two calls were needed because MediaWiki's plain-text endpoint produces well-structured prose but strips out infobox tables, while the HTML endpoint preserves the infobox but yields lower-quality body text. The HTML was parsed with BeautifulSoup to extract the infobox table (aliases, allies, rivals, years active, etc.), which was appended to the article text. Each article's content and versioned URL were saved to disk.

\subsection{Network extraction via LLM}
Each article was sent to the DeepSeek API with a structured prompt that enforced a fixed output schema. For each article, the LLM extracted:
\begin{itemize}[leftmargin=2em]
  \item \textbf{Entities (nodes):} standardized name, aliases, type, contextual description, and time period. The type was constrained to one of nine canonical values (cartel, mafia, gang, motorcycle club, clan, triad, militia, faction, terrorist organization) or ``other''.
  \item \textbf{Relationships (edges):} source organization, target organization, relationship type (alliance, rivalry, or other), optional detail, context, and time period. The prompt explicitly mapped cooperation to alliance and conflict to rivalry, reserving ``other'' when the relationship is genuinely neither cooperation nor conflict.
  
\end{itemize}
Articles longer than 2{,}500 words were chunked, with the article's opening paragraph passed as context to every chunk to preserve entity resolution across chunks. Extraction ran in parallel across 50 concurrent workers and returned a JSON object of nodes and edges for each article.

\subsection{Cleanup pipeline}

All JSON files were merged into a single global network. Nodes were deduplicated by name (case-insensitive), with aliases, descriptions, and source article references merged across duplicates. Edges were deduplicated by (source, target, relationship type, detail). Every node and edge retains the Wikipedia URL it was extracted from. A curated dictionary maps known variant spellings to their canonical form (e.g., ``Medellin Cartel'' $\rightarrow$ ``Medell\'in Cartel'', ``Hells Angels MC'' $\rightarrow$ ``Hells Angels'', ``FARC-EP'' $\rightarrow$ ``FARC'').

Because the prompt enforced canonical types and relationship classifications directly, the cleanup pipeline focuses on work the LLM cannot do per-article:

\begin{itemize}[leftmargin=2em]
  \item \textbf{Organization type overrides.} Organizations the LLM mistyped had their type forced to the correct value via a hand-curated mapping.
  \item \textbf{Hand-curated exclusions.} Non-criminal entities that occasionally slipped through extraction (governments, NGOs, political parties) were removed via a curated exclusion list. The same list was used to remove defunct organizations: those with documented dissolution dates and no evidence of continued activity.
  \item \textbf{Generic node filtering.} Umbrella terms referring to categories rather than specific organizations (e.g., ``Russian organized crime'', ``Colombian drug cartels'') were removed. Specific organizations with similar names (e.g., ``Mexican Mafia'', ``Albanian Mafia'') were preserved.
  \item \textbf{Type-name safety net.} A small fallback map catches occasional non-canonical type names that the LLM emits despite the prompt constraints (e.g., \texttt{crime\_family} $\rightarrow$ \texttt{mafia}). 
   \item \textbf{Chronological edge deduplication.} Edges connecting the same organizations were deduplicated. When several articles described the same pair, or when the relationship between two organizations changed over time, multiple edges accumulate. Therefore, I pass it through the llm: if every edge has a parseable date, keep the one with the latest end date. If some edges have dates and others don't, ask the LLM to pick the most current. If no edges have dates, ask the LLM to pick the most current.
\end{itemize}

\subsection{Dataset overview}

The final dataset contains 771 source articles with full Wikipedia URL provenance, from which I extracted:

\begin{itemize}[leftmargin=2em]
  \item \textbf{1{,}890 entities} across nine organization types: gangs (39.7\%), clans (13.5\%), mafias (13.4\%), motorcycle clubs (10.4\%), factions (9.3\%), cartels (4.3\%), militias (4.2\%), triads (2.8\%), terrorist organizations (2.2\%).
  \item \textbf{3{,}354 relationships:} 2{,}293 alliances (68\%) and 1{,}061 rivalries (32\%). All edges are undirected.
\end{itemize}

\newpage

% ═══════════════════════════════════════════════════════
\section{Network structure}\label{sec:network}
% ═══════════════════════════════════════════════════════

\subsection{The full network}\label{sec:layers}

Alliance and rivalry are very different types of relationships. While the rest of this analysis treats them as two separate networks, seeing them combined first gives a helpful overview of the overall scale of these organizations and their connections.

It is important to note that this network reflects the exact state of these organizations right now. It shows the active landscape of global organized crime as of April 27, 2026.

\begin{figure}[H]
  \centering
  \includegraphics[width=\textwidth]{images/crimenet_both_full.pdf}
  \caption{Complete CRIMENET with both relationship types. Nodes are criminal organizations, sized by betweenness centrality. Blue edges are alliances, red edges are rivalries. The entire network contains 1{,}890 nodes and 3{,}354 edges.}
  \label{fig:multiplex}
\end{figure}

By separating the network, we can see what cooperation and conflict look like independently. Figure \ref{fig:layers} shows the alliance and rivalry layers side by side.

\begin{figure}[H]
  \centering
  \includegraphics[width=\textwidth]{images/join_networks.png}
  \caption{Alliance (left, blue) and rivalry (right, red) networks shown separately.}
  \label{fig:layers}
\end{figure}

The first question is a basic one: how do the two networks differ in size, reach, and shape? Table~\ref{tab:structural} reports the structural properties.

\begin{table}[H]
\centering
\caption{Structural properties of the alliance and rivalry networks.}
\label{tab:structural}
\small
\begin{tabular}{@{}lcc@{}}
\toprule
\textbf{Metric} & \textbf{Alliance} & \textbf{Rivalry} \\
\midrule
Active nodes (deg $>$ 0) & 1{,}517 & 876 \\
Edges & 2{,}293 & 1{,}061 \\
Density & 0.00199 & 0.00277 \\
Mean degree & 3.02 & 2.42 \\
Median degree & 1 & 1 \\
Max degree & 134 & 53 \\
Components & 68 & 102 \\
LCC nodes & 1{,}292 (85.2\%) & 478 (54.6\%) \\
LCC edges & 2{,}130 (92.9\%) & 744 (70.1\%) \\
LCC density & 0.00255 & 0.00653 \\
Pseudo-diameter (LCC) & 12 & 13 \\
Global clustering & 0.0906 & 0.0859 \\
Mean local clustering & 0.2957 & 0.1893 \\
Degree assortativity & $-$0.1236 & $-$0.0796 \\
Max $k$-core & 8 & 4 \\
Max betweenness & 0.187749 & 0.134382 \\
Max PageRank & 0.027264 & 0.019312 \\
\bottomrule
\end{tabular}
\end{table}

The alliance network has more active nodes (1{,}517 vs 876) and edges (2{,}293 vs 1{,}061) and its LCC covers 85\% of active nodes versus 55\%. Cooperation reaches more of the network than conflict does. Global clustering is 0.0906 vs 0.0859, mean local clustering 0.2957 vs 0.1893. Pseudo-diameters are close (12 vs 13). One detail worth flagging: the rivalry LCC is denser than the alliance LCC (0.00653 vs 0.00255), even though the rivalry network is sparser overall. Conflict packs into smaller, denser clusters; cooperation spreads thinner across a larger one.

Beyond the overall shape of the networks, it is also useful to measure how often similar groups interact. Type assortativity (Newman's categorical \href{https://journals.aps.org/pre/abstract/10.1103/PhysRevE.67.026126}{$r$}) is higher for rivalries than alliances ($r = 0.64$ vs $0.49$, $p < 0.001$, edge bootstrap). Organizations fight within their type more often than they ally within it.

\subsection{Network centrality}

To identify the most prominent organizations in each network, I ranked them using three standard metrics: degree, betweenness, and PageRank. Figure~\ref{fig:top10} shows the top ten organizations in each network for all three metrics.  

\begin{figure}[H]
  \centering
  \includegraphics[width=\textwidth]{images/crimenet_top10_centrality.pdf}
  \caption{Top 10 organizations by degree (top), betweenness centrality (middle), and PageRank (bottom) for the alliance (left) and rivalry (right) networks. Node colors indicate organizational type.}
  \label{fig:top10}
\end{figure}

Hells Angels top every ranking in both networks; no other organization does.
 
Alliance: Hells Angels, 'Ndrangheta, and Sinaloa Cartel lead, with the rest of the top slots held by mafias (Camorra, Gambino crime family, American Mafia, Cosa Nostra, Mexican Mafia, Bonanno crime family, Chicago Outfit) and the occasional cartel. Only the Crips and Bloods break into the alliance betweenness top ten, and no gang reaches the alliance PageRank top ten.
 
Rivalry: motorcycle clubs move to the front (Bandidos 2nd in degree and PageRank; Outlaws across rankings). Gangs fill much of the list (Crips, Bloods, Latin Kings, Mara Salvatrucha, Florencia 13, Norteños, Tiny Rascal Gang, Maniac Latin Disciples). Cartels stay prominent (Jalisco New Generation Cartel, Sinaloa Cartel). It is worth noticing the only mafia present is Mexican Mafia, high in rivalry brokerage (2nd in betweenness) and absent from the alliance top ten.

Alliance and rivalry brokerage are held by mostly different organizations. Since gangs and motorcycle clubs clearly dominate the top of the rivalry network, the next step is to aggregate these metrics to see how each organizational type performs as a whole. The top-ten lists are dominated by a few organizations. Do the patterns hold when averaged across all organizations of a type? Figure~\ref{fig:type_centrality} shows the mean of each metric by type.

\begin{figure}[H]
  \centering
  \includegraphics[width=\textwidth]{images/crimenet_importance_by_type.pdf}
  \caption{Mean degree (top), betweenness (middle), and PageRank (bottom) by organizational type, with standard error of the mean, for the alliance (left) and rivalry (right) networks.}
  \label{fig:type_centrality}
\end{figure}

By type, mafias lead alliance mean degree and PageRank, while cartels lead alliance mean betweenness; mafias come second on betweenness, with triads and terrorist organizations rounding out the upper third. Factions sit at the bottom of mean alliance betweenness, with gangs in the middle of the pack. On the rivalry side, cartels lead in mean degree and mean PageRank, motorcycle clubs lead in mean betweenness, and mafias are third in PageRank but only fifth in degree and betweenness. The same handful of types dominates both layers, but the ordering depends on which relationship is being measured.

Do prestige and brokerage identify the same organizations, or different ones? Figure~\ref{fig:pr_bc_scatter} plots PageRank against betweenness for every active organization.

\begin{figure}[H]
  \centering
  \includegraphics[width=\textwidth]{images/crimenet_scatter_combined.pdf}
  \caption{PageRank versus betweenness centrality for active organizations in the alliance (left) and rivalry (right) networks, colored by type.}
  \label{fig:pr_bc_scatter}
\end{figure}

In the alliance top-right corner: Hells Angels, Sinaloa Cartel, 'Ndrangheta, Gambino crime family, Cosa Nostra, American Mafia, Mexican Mafia. The same organizations score high on both axes. In the rivalry plot, Hells Angels sits alone at the peak; high-betweenness organizations (Bloods, Mara Salvatrucha, Mongols, Mexican Mafia) have middling PageRank, and high-PageRank organizations (Bandidos, Jalisco New Generation Cartel, Sinaloa Cartel, Crips) have middling betweenness. Rivalry prestige and rivalry brokerage are held by different organizations.  Again, we observe that alliance is more type-mixed wheares rivalry is more dominated by gangs (blue circles) and motorcycle clubs. With the exception of Mexican Mafia. 

Now let's take a look at the core of the networks.

\newpage

\subsection{$k$-core decomposition}

Which types sit at the core of each network, and which at the periphery? The \href{https://www.sciencedirect.com/science/article/abs/pii/037887338390028X}{$k$-core} of a network is the largest subgraph in which every vertex has at least $k$ neighbors. $k$-core decomposition applies this recursively, peeling off lower-connectivity vertices to reveal a nested shell structure. Figure~\ref{fig:kcore} shows the type composition of each shell.

\begin{figure}[H]
  \centering
  \includegraphics[width=\textwidth]{images/crimenet_kcore_composition.pdf}
  \caption{Type composition of $k$-core shells for the alliance (left) and rivalry (right) networks. Numbers above bars indicate shell size.}
  \label{fig:kcore}
\end{figure}

The alliance network reaches $k=8$; the rivalry network reaches $k=4$. The alliance core runs deeper. Gangs dominate all the cores of the rivalry network. On the other hand, the deepest alliance shell ($k=8$) contains 22 organizations, almost all mafias---including the major American Mafia crime families (Gambino crime family, Genovese crime family, Bonanno crime family, Lucchese crime family, Colombo crime family), other crime families and, of course, Hells Angels. The decomposition surfaces the Commission structure without supervision. The rivalry $k=4$ shell contains 43 organizations, predominantly gangs. The rivalry core is shallower and more type-homogeneous shell by shell, consistent with the higher type assortativity reported in Section~\ref{sec:layers}.

\subsection{Signed pairs}\label{sec:signed}

So far I've looked at the two networks separately. Putting them side by side raises a different question: which organizations are pure cooperators, which are pure antagonists, and which keep a mixed book? I collapse parallel edges at the organization-pair level, giving each pair a signed label (ally, rival, or both). 2{,}293 pairs have an alliance edge, 1{,}061 have a rivalry edge, 0 have both. Table~\ref{tab:signed} lists the 15 organizations with the highest total signed degree (allies $+$ rivals): the most involved, regardless of direction. 

THIS IS IMPORTANT: here add a little comment that this is only for this gathered dataset, they might have more (or any) allies/rivals that the data didn't compile it.

\begin{table}[H]
\centering
\caption{Top 15 organizations by total signed degree (allies $+$ rivals).}
\label{tab:signed}
\small
\begin{tabular}{@{}llcc@{}}
\toprule
\textbf{Name} & \textbf{Type} & \textbf{Allies} & \textbf{Rivals} \\
\midrule
Hells Angels & motorcycle club & 134 & 53 \\
Sinaloa Cartel & cartel & 66 & 18 \\
Bandidos & motorcycle club & 33 & 38 \\
Outlaws & motorcycle club & 33 & 27 \\
Crips & gang & 30 & 25 \\
Mexican Mafia & mafia & 40 & 14 \\
Bloods & gang & 32 & 18 \\
Jalisco New Generation Cartel & cartel & 21 & 23 \\
Latin Kings & gang & 24 & 19 \\
Gulf Cartel & cartel & 30 & 13 \\
Los Zetas & cartel & 28 & 12 \\
Sureños & gang & 29 & 8 \\
Yamaguchi-gumi & mafia & 20 & 12 \\
Primeiro Comando da Capital & faction & 18 & 12 \\
Mara Salvatrucha & gang & 13 & 16 \\
\bottomrule
\end{tabular}
\end{table}

Pure cooperators (only allies, 0 rivals) are almost entirely mafias and clans. Comment result below:

% Top 10 PURE COOPERATORS (5+ allies, 0 rivals):
%   Cosa Nostra                                    + 42 / -  0  type=mafia
%   Barbaro 'ndrina                                + 22 / -  0  type=mafia
%   Detroit Partnership                            + 19 / -  0  type=mafia
%   Patriarca crime family                         + 18 / -  0  type=mafia
%   Corsican mafia                                 + 17 / -  0  type=mafia
%   Jamaican posse                                 + 17 / -  0  type=gang
%   Queen City Kings                               + 17 / -  0  type=gang
%   New Orleans crime family                       + 14 / -  0  type=mafia
%   Trafficante crime family                       + 14 / -  0  type=mafia
%   Five Families                                  + 13 / -  0  type=mafia

We have only 2 pure antagonists (only rivals, 0 allies) and both are gangs. Comment result below:
% Top 10 PURE ANTAGONISTS (5+ rivals, 0 allies):
%   Fresno Bulldogs                                +  0 / -  6  type=gang
%   Trinitarios                                    +  0 / -  6  type=gang

\newpage

\subsection{Role persistence across networks}

If an organization is a broker in the alliance network, is it also a broker in the rivalry network? To find out, I rank each organization active in both layers by betweenness within each network and compare the ranks. Perfect transfer would put every point on the diagonal in Figure~\ref{fig:role_persistence}.

\begin{figure}[H]
  \centering
  \includegraphics[width=\textwidth]{images/crimenet_role_persistence.pdf}
  \caption{Betweenness rank in the alliance network ($x$-axis) versus rivalry network ($y$-axis) for the 376 organizations with nonzero betweenness in at least one layer. Colors indicate organizational type; node size is proportional to degree.}
  \label{fig:role_persistence}
\end{figure}

We can clearly see that most points don't fall on the diagonal. I score persistence as $\max(\text{rank}_a, \text{rank}_r)$, so an organization gets a low score only if it is highly ranked in both layers, with ties broken by $|\Delta|$. Table~\ref{tab:persistent} lists the top 10.

\begin{table}[H]
\centering
\caption{Top 10 most persistent organizations by betweenness centrality (highly ranked in both networks).}
\label{tab:persistent}
\small
\begin{tabular}{@{}llccc@{}}
\toprule
\textbf{Organization} & \textbf{Type} & \textbf{Alliance Rank} & \textbf{Rivalry Rank} & $\boldsymbol{|\Delta|}$ \\
\midrule
Hells Angels & Motorcycle Club & 1 & 1 & 0 \\
Crips & Gang & 8 & 4 & 4 \\
Bloods & Gang & 9 & 5 & 4 \\
Mexican Mafia & Mafia & 11 & 2 & 9 \\
Bandidos & Motorcycle Club & 12 & 7 & 5 \\
Outlaws & Motorcycle Club & 13 & 11 & 2 \\
Sinaloa Cartel & Cartel & 3 & 13 & 10 \\
Latin Kings & Gang & 15 & 9 & 6 \\
Los Zetas & Cartel & 17 & 14 & 3 \\
14K Triad & Triad & 14 & 19 & 5 \\
\bottomrule
\end{tabular}
\end{table}

The clearest exception is Hells Angels: rank 1 on both axes, and the system's outlier on every measure I've looked at (top of every centrality ranking in both networks, top of total signed degree).

Alternatively, the points that fall \emph{off} the diagonal split into two groups: alliance-dominant brokers (high alliance rank, low rivalry rank) and rivalry-dominant brokers (the reverse). Tables~\ref{tab:alliance_dominant} and~\ref{tab:rivalry_dominant} list the top 10 in each direction.

\begin{table}[H]
\centering
\caption{Top 10 alliance-dominant organizations by betweenness rank divergence.}
\label{tab:alliance_dominant}
\small
\begin{tabular}{@{}llccc@{}}
\toprule
\textbf{Organization} & \textbf{Type} & \textbf{Alliance Rank} & \textbf{Rivalry Rank} & $\boldsymbol{\Delta}$ \\
\midrule
Serbian mafia & Mafia & 38 & 359 & $-$321 \\
Secondigliano Alliance & Clan & 51 & 358 & $-$307 \\
The British & Gang & 66 & 367 & $-$301 \\
Nazi Lowriders & Gang & 42 & 339 & $-$297 \\
Hard Livings & Gang & 16 & 310 & $-$294 \\
Greek mafia & Mafia & 20 & 304 & $-$284 \\
Piru Street Boys & Gang & 65 & 348 & $-$283 \\
Shining Path & Terrorist Org. & 77 & 360 & $-$283 \\
Hidden Valley Kings & Gang & 29 & 311 & $-$282 \\
Lucchese crime family & Mafia & 47 & 329 & $-$282 \\
\bottomrule
\end{tabular}
\end{table}

Alliance-dominant brokers (Table~\ref{tab:alliance_dominant}): discuss table 

\begin{table}[H]
\centering
\caption{Top 10 rivalry-dominant organizations by betweenness rank divergence.}
\label{tab:rivalry_dominant}
\small
\begin{tabular}{@{}llccc@{}}
\toprule
\textbf{Organization} & \textbf{Type} & \textbf{Alliance Rank} & \textbf{Rivalry Rank} & $\boldsymbol{\Delta}$ \\
\midrule
Toonerville Rifa 13 & Gang & 372 & 26 & 346 \\
West Side Piru & Gang & 374 & 30 & 344 \\
Grape Street Watts Crips & Gang & 342 & 10 & 332 \\
Santa Monica 13 & Gang & 365 & 70 & 295 \\
Los Chone Killers & Gang & 353 & 63 & 290 \\
Playboys & Gang & 359 & 77 & 282 \\
Logan Heights Gang & Gang & 352 & 74 & 278 \\
La Unión Tepito & Cartel & 350 & 80 & 270 \\
Original Red Devils MC & Motorcycle Club & 299 & 33 & 266 \\
Tango Blast & Gang & 370 & 105 & 265 \\
\bottomrule
\end{tabular}
\end{table}

Rivalry-dominant brokers (Table~\ref{tab:rivalry_dominant}): discuss table (also saying that mostly are gangs)

\subsection{Edge composition by type}

The centrality and $k$-core analyses treat each organization individually. A complementary question: of all edges incident on organizations of a given type, what fraction are alliances? This measures how a type spends its relational footprint as a whole. For each type, I pool all edges incident on organizations of that type and compute the fraction that are alliances. Confidence intervals are bootstrap 95\% intervals obtained by resampling organizations within each type with replacement ($B = 2000$).

\begin{figure}[H]
  \centering
  \includegraphics[width=0.8\textwidth]{images/crimenet_alliance_share_by_type.pdf}
  \caption{Alliance share of edges by organizational type, edge-weighted (each edge counts equally). Error bars are bootstrap 95\% confidence intervals over organizations ($B = 2000$). The dashed line marks 50\%.}
  \label{fig:edge_composition}
\end{figure}

comment results below:

% Type                          n   Ally %             95% CI
% ------------------------------------------------------------
% Mafia                       263    86.1%       [82.5, 89.1]
% Terrorist Organization       28    82.5%       [71.2, 90.1]
% Faction                     192    79.1%       [73.5, 85.0]
% Clan                        251    77.9%       [73.6, 82.1]
% Militia                      46    71.1%       [62.4, 80.3]
% Triad                        23    69.2%       [54.7, 77.5]
% Cartel                       88    66.2%       [59.6, 72.0]
% Motorcycle Club             233    63.6%       [57.7, 68.3]
% Gang                        766    54.8%       [52.1, 57.5]


% ═══════════════════════════════════════════════════════
\section{Limitations}\label{sec:limitations}
% ═══════════════════════════════════════════════════════

\begin{enumerate}[leftmargin=2em]
  \item \textbf{Wikipedia coverage bias.} Wikipedia skews English-language and Western-centric. Well-documented organizations appear more connected. Coverage gaps likely exist in under-reported regions.
  \item \textbf{LLM extraction errors.} The sanitization pipeline catches obvious problems, but I have not run a systematic precision/recall evaluation against expert-labeled data.
  \item \textbf{Temporal conflation.} The network aggregates relationships across all time periods. An alliance from the 1970s and one from 2020 are treated identically.
  \item \textbf{Type classification ambiguity.} The boundary between ``cartel'' and ``militia'', or ``gang'' and ``faction'', is often fuzzy.
\end{enumerate}

% ═══════════════════════════════════════════════════════
\section{Conclusion}\label{sec:conclusion}
% ═══════════════════════════════════════════════════════

The alliance network is denser and has a deeper $k$-core ($k=11$ vs $k=7$). The rivalry network is sparser and more fragmented, but its connected portion is denser. Type assortativity is higher in the rivalry network ($r = 0.62$ vs $0.52$).

Mafias lead the alliance network on every centrality metric, occupy the deepest $k$-core shells, and 86\% of their edges are alliances. Gangs hold most rivalry-dominant brokerage positions, and 45\% of their edges are rivalries. Brokerage roles are mostly layer-specific: alliance-dominant brokers are mostly mafias, rivalry-dominant brokers are mostly gangs. A smaller group of organizations brokers in both layers, and that group is type-mixed rather than type-clustered. Hells Angels are the single most central organization across the entire dataset.

% \bibliographystyle{unsrt}
% \bibliography{refs}

\end{document}